\songtitle{La Polko de Eva (Ievan Polkka)}{Eino Kettunen}{1928}
{\color{gray}Traduko de Joop Kiefte}

\medskip

\begin{paracol}{2}
\begin{guitar}
	\songsection{Strofo 1}
  Kiam sonas tre ĝoje la polka sonado
  la piedoj ekmovas tiel ĉi.
  Kaj Eva ja estis sub garda patrino
  sed Eva al ŝi bone kaŝis sin.
  Ĉar kiu aŭskultas al sia patrino
  se tiu tuj diras ne al la festo.
  Mi volas danci kaj iri nun
  al fekmojosa dancfesten'

	\songsection{Strofo 2}
  Jen Eva ridadis, muziko ludaĉis
  kaj homa amaso kolektas sin.
  Ĉiuj bonegis, ja ĉiuj belegis
  sed unu belulo kaptis ŝin.
  Ĉar kiu ne volas peni treege,
  kiu ne volas danci freneze
  se tiu belulo tiom belas
  en tiu ĉi dancfesten'?

	\songsection{Strofo 3}
  Panjo, malĝoja, eniris dormĉambron
  kaj iel ploretis kaj kantis, nu.
  Tiel la ulo ricevis la spacon
  en najbara domo por amplezur'
  Ĉar kiu ja sekvas gepatrajn konsilojn,
  kiu aŭskultas riproĉajn avertojn
  se tiu belulo tiom belas
  en tiu ĉi dancfesten'?
\end{guitar}

\switchcolumn

\begin{guitar}
	\songsection{Strofo 4}
  Kiam haltis muziko estis nur komenco
  la knabo ekumis, kial ne?
  Kiam li ŝin post la festo kunportis
  patrino atendis kolerege.
  Sed mi Evan diris, ne malĝojiĝu,
  mi ja revenos, min ne forgesu,
  se tia belulo tiom belas
  en tia ĉi dancfesten'?

	\songsection{Strofo 5}
  Mi al patrino avertis nun ĉesi
  por ke mi atentu pri mia far'
  Vi nun ekiru, en ĉambro nur resti
  por ke estu sekura la aferar'.
  Ĉar tiu ĉi knabo estas sovaĝulo
  kiu ja faras kion li volas,
  ja tia belulo tiom belas
  en tia ĉi dancfesten'!

	\songsection{Strofo 6}
  Nur afereton al vi mi postlasas:
  vi ja ne facile kaptos min.
  Iru al ajna volata direkto
  sed Eva kaj mi kombinos nin.
  Ĉar tiu ĉi knabo ne estas freneza
  tiu ĉi knabo plenumas promesojn
  ja tiu belulo tiom belis
  en tiu ĉi dancfesten'.
\end{guitar}
\end{paracol}
